\title{Set-based Minesweeper}
\date{}
\author{peppermintpatty5}

\documentclass[12pt]{article}
\setlength\parindent{0pt}

\usepackage{amsfonts}


\begin{document}

\maketitle
\tableofcontents

\section{Data Structures}

A minesweeper grid is defined by three sets of pairs of integers $R, F, M \subset \mathbb{Z} ^ 2$ which denote the indices of revealed cells, flagged cells, and cells that are mines respectively. It is always the case that $R \cap F = \emptyset$ because a cell cannot be both revealed and flagged.

\section{Algorithms}

Let $A$ be the set of cells adjacent to $x = (r, c)$:

$$ A = \left ( \left \{ r - 1, r, r + 1 \right \} \times \left \{ c - 1, c, c + 1 \right \} \right ) \setminus \left \{ \left ( r, c \right ) \right \} $$

\subsection{Reveal}

If $x$ is not flagged:

$$ x \notin F $$

Then add $x$ to $R$:

$$ R \cup \left \{ x \right \} $$

\subsection{Flag}

If $x$ has not been revealed:

$$ x \notin R $$

Then if $x$ is flagged, then it will be un-flagged and vice versa. This is done rather neatly by updating $F$ via symmetric difference:

$$ F \triangle \left \{ x \right \} $$

\subsection{Chord}

If $x$ has been revealed and is not a mine and the number of adjacent flags plus the number of adjacent revealed mines equals the number of adjacent mines:

$$ x \in R \wedge x \notin M \wedge \left | A \cap F \right | + \left | A \cap R \cap M \right | = \left | A \cap M \right | $$

Then add the adjacent non-flagged cells to $R$:

$$ R \cup \left ( A \setminus F \right ) $$

\end{document}
